\documentclass[12pt]{article}
\usepackage{fullpage,amsmath,amssymb,amsthm,hyperref,amsfonts}
\newtheorem{theorem}{Theorem}
\newtheorem{lemma}[theorem]{Lemma}
\newtheorem{proposition}[theorem]{Proposition}
\theoremstyle{definition}
\newtheorem{definition}[theorem]{Definition}
\newtheorem{remark}[theorem]{Remark}

\title{Equivalence of the $\Phi^4_3$ Measure under Smooth Shifts}
\author{}
\date{}

\begin{document}
\maketitle

\begin{abstract}
Let $\mathbb{T}^3$ be the three-dimensional unit torus, $\mu$ the $\Phi^4_3$ measure on
$\mathcal{D}'(\mathbb{T}^3)$, and $\psi\colon\mathbb{T}^3\to\mathbb{R}$ a nonzero smooth function.
The shift map $T_\psi(u)=u+\psi$ induces the pushforward $(T_\psi)^*\mu$.
We prove that $\mu$ and $(T_\psi)^*\mu$ are equivalent measures, i.e., they have identical null
sets.
\end{abstract}

\section{Setup and Definitions}

\subsection{The reference Gaussian free field}

Fix $m>0$.  The \emph{Gaussian free field} (GFF) $\mu_0$ on $\mathcal{D}'(\mathbb{T}^3)$ is the
centred Gaussian measure with covariance operator
\[
  C \;:=\; (-\Delta + m^2)^{-1} \;:\; L^2(\mathbb{T}^3)\to L^2(\mathbb{T}^3),
  \label{eq:covariance}
\]
where $\Delta$ denotes the Laplace--Beltrami operator on $\mathbb{T}^3$.  Concretely, if
$\{e_k\}_{k\in\mathbb{Z}^3}$ are the standard Fourier characters on $\mathbb{T}^3$ and
$\xi_k\sim N(0,1)$ are i.i.d., then $\phi = \sum_k \xi_k (|k|^2+m^2)^{-1/2} e_k$ realizes $\mu_0$.
The operator $C$ is a bounded, positive, self-adjoint operator on $L^2(\mathbb{T}^3)$ with
spectrum $(0, m^{-2}]$.

\subsection{The Cameron--Martin space of $\mu_0$}

\begin{definition}[Cameron--Martin space]
\label{def:CM}
The \emph{Cameron--Martin space} of a Gaussian measure $\nu$ with covariance operator $C$
on a Hilbert space $\mathcal{X}$ is
\[
  \mathcal{H}(\nu) \;:=\; C^{1/2}(\mathcal{X}),
\]
equipped with the inner product
$\langle h_1, h_2\rangle_{\mathcal{H}} = \langle C^{-1/2}h_1,\, C^{-1/2}h_2\rangle_{\mathcal{X}}$.
\end{definition}

For the GFF $\mu_0$ with $C=(-\Delta+m^2)^{-1}$:
\[
  \mathcal{H}(\mu_0) \;=\; (-\Delta+m^2)^{-1/2}(L^2(\mathbb{T}^3))
  \;=\; H^1(\mathbb{T}^3),
\]
with squared norm
\begin{equation}
  \|h\|_{\mathcal{H}}^2
  \;=\; \langle C^{-1}h,\, h\rangle_{L^2}
  \;=\; \langle (-\Delta+m^2)h,\, h\rangle_{L^2}
  \;=\; \|\nabla h\|_{L^2}^2 + m^2\|h\|_{L^2}^2
  \;=\; \|h\|_{H^1_m}^2.
  \label{eq:CMnorm}
\end{equation}
The last equality in \eqref{eq:CMnorm} uses the identity
$\langle -\Delta h, h\rangle_{L^2} = \|\nabla h\|_{L^2}^2$ (integration by parts on $\mathbb{T}^3$).

Since $\psi\in C^\infty(\mathbb{T}^3)$, we have $\psi\in H^s(\mathbb{T}^3)$ for every $s\geq 0$, and in
particular $\psi\in H^1(\mathbb{T}^3) = \mathcal{H}(\mu_0)$.

\subsection{The $\Phi^4_3$ measure}

The \emph{$\Phi^4_3$ measure} $\mu$ is the probability measure on $\mathcal{D}'(\mathbb{T}^3)$
formally written
\[
  d\mu(\phi) \;=\; Z^{-1} e^{-V(\phi)}\, d\mu_0(\phi),
\]
where $V(\phi)$ is the Wick-renormalized quartic interaction
\[
  V(\phi) \;=\; \lambda \int_{\mathbb{T}^3} {:\phi(x)^4:}_{\mu_0}\, dx
  \;+\; a_2 \int_{\mathbb{T}^3} {:\phi(x)^2:}_{\mu_0}\, dx
  \;+\; a_0,
  \label{eq:interaction}
\]
with coupling constant $\lambda>0$, renormalization constants $a_2, a_0\in\mathbb{R}$ chosen to
cancel the ultraviolet divergences in $d=3$, and $Z = \int e^{-V(\phi)}\,d\mu_0(\phi)$
the normalizing partition function.  Here ${:\phi(x)^k:}_{\mu_0}$ denotes the $k$-th Wick power
of $\phi$ with respect to the covariance $C$ of $\mu_0$.

The rigorous construction of $\mu$ (showing that $Z\in(0,\infty)$ and that $e^{-V}\in L^1(\mu_0)$)
has been carried out by Nelson \cite{Nelson1966}, by Glimm and Jaffe
\cite{GlimmJaffe1987} via the Euclidean approach, and more recently via Hairer's regularity
structures \cite{Hairer2014} and via paracontrolled distributions
\cite{GubinelliHofmanova2019}.  The measure $\mu$ is supported on $\mathcal{C}^{-1/2-\varepsilon}(\mathbb{T}^3)$
for every $\varepsilon>0$, in the sense that $\mu$-almost every $\phi$ belongs to
$\mathcal{C}^{-1/2-\varepsilon}(\mathbb{T}^3)$.

\section{Main Result and Proof}

\begin{theorem}[Equivalence under smooth shift]
\label{thm:main}
Let $\psi\in C^\infty(\mathbb{T}^3)$ be nonzero.  Then $\mu$ and $(T_\psi)^*\mu$ are equivalent
probability measures on $\mathcal{D}'(\mathbb{T}^3)$.  Explicitly, $(T_\psi)^*\mu \ll \mu$ and
$\mu \ll (T_\psi)^*\mu$, with Radon--Nikodym derivative
\begin{equation}
  \frac{d(T_\psi)^*\mu}{d\mu}(\phi)
  \;=\; R_\psi(\phi)\cdot e^{W(\phi-\psi;\,\psi)},
  \label{eq:RNderiv}
\end{equation}
where
\begin{align}
  R_\psi(\phi)
  &\;:=\; \exp\!\Bigl(\langle (-\Delta+m^2)\psi,\,\phi\rangle_{L^2}
        - \tfrac{1}{2}\langle (-\Delta+m^2)\psi,\,\psi\rangle_{L^2}\Bigr),
  \label{eq:Rpsi} \\
  W(\xi;\,\psi)
  &\;:=\; V(\xi+\psi) - V(\xi).
  \label{eq:Wpsi}
\end{align}
\end{theorem}

The proof proceeds through three lemmas.

\begin{lemma}[Cameron--Martin for the GFF]
\label{lem:CM}
Let $\psi\in H^1(\mathbb{T}^3)$.  Then $(T_\psi)^*\mu_0 \sim \mu_0$, with
\begin{equation}
  \frac{d(T_\psi)^*\mu_0}{d\mu_0}(\phi)
  \;=\; R_\psi(\phi),
  \label{eq:CMLGFF}
\end{equation}
where $R_\psi$ is given by \eqref{eq:Rpsi}.  In particular, $R_\psi(\phi) > 0$ for all
$\phi\in\mathcal{D}'(\mathbb{T}^3)$.
\end{lemma}

\begin{proof}
We apply the Cameron--Martin theorem in the following form (see Bogachev
\cite[Theorem 2.4.5 and Corollary 2.4.3]{Bogachev1998}):

\textbf{Cameron--Martin theorem.}
\emph{Let $\nu$ be a centred Gaussian measure on a locally convex space $\mathcal{X}$ with
Cameron--Martin space $\mathcal{H}(\nu)$.  For $h\in\mathcal{H}(\nu)$, the shifted measure
$(T_h)^*\nu$ is equivalent to $\nu$, with Radon--Nikodym derivative}
\[
  \frac{d(T_h)^*\nu}{d\nu}(\phi)
  \;=\; \exp\!\bigl(\hat{h}(\phi) - \tfrac{1}{2}\|h\|_{\mathcal{H}}^2\bigr),
\]
\emph{where $\hat{h}\in L^2(\nu)$ is the measurable linear functional associated to $h$ via the
isometric embedding $\mathcal{H}(\nu)\hookrightarrow L^2(\nu)$, defined by
$\hat{h}(\phi) := \langle C^{-1}h,\,\phi\rangle$ (interpreted in the sense of the Cameron--Martin
pairing).}

We verify the hypotheses in our setting.

\textit{Hypothesis 1: $\nu = \mu_0$ is a centred Gaussian measure.}  By definition, $\mu_0$ is the
centred Gaussian with covariance $C = (-\Delta+m^2)^{-1}$.  This is a well-defined Gaussian
measure on $\mathcal{D}'(\mathbb{T}^3)$ since $C$ is a positive trace-class operator on
$L^2(\mathbb{T}^3)$ (the eigenvalues $(|k|^2+m^2)^{-1}$ are summable in 3 dimensions:
$\sum_{k\in\mathbb{Z}^3}(|k|^2+m^2)^{-1}$ converges, cf.\ e.g.\ \cite[p.~25]{GlimmJaffe1987}).

\textit{Hypothesis 2: $\psi\in\mathcal{H}(\mu_0) = H^1(\mathbb{T}^3)$.}
We have $\psi\in C^\infty(\mathbb{T}^3)\subset H^1(\mathbb{T}^3) = \mathcal{H}(\mu_0)$, as shown in
Section~1.2.

\textit{Identification of $\hat\psi$.}  For $h = \psi\in H^1(\mathbb{T}^3) = C^{1/2}(L^2)$, the
element $C^{-1}h = (-\Delta+m^2)\psi\in L^2(\mathbb{T}^3)$ (using $\psi\in C^\infty\subset H^2$),
and the functional $\hat\psi(\phi) = \langle C^{-1}\psi,\,\phi\rangle_{L^2} = \langle (-\Delta+m^2)\psi,\,\phi\rangle_{L^2}$.
Since $(-\Delta+m^2)\psi\in L^2(\mathbb{T}^3)\subset L^2(\mathbb{T}^3)$, this pairing is
well-defined $\mu_0$-a.s.~(as $L^2$-functions act continuously on distributions in
$\mathcal{D}'$).

\textit{Identification of $\|\psi\|_{\mathcal{H}}^2$.}
By \eqref{eq:CMnorm},
$\|\psi\|_{\mathcal{H}}^2 = \langle C^{-1}\psi,\,\psi\rangle_{L^2} = \langle (-\Delta+m^2)\psi,\,\psi\rangle_{L^2}$.

Applying the Cameron--Martin theorem, $(T_\psi)^*\mu_0\sim\mu_0$ with the density
\eqref{eq:CMLGFF}.  Strict positivity $R_\psi(\phi)>0$ holds for all $\phi$ because $R_\psi$
is an exponential of a real number.
\end{proof}

\begin{lemma}[Absolute continuity of densities w.r.t.\ $\mu_0$]
\label{lem:densities}
Both $\mu$ and $(T_\psi)^*\mu$ are absolutely continuous with respect to $\mu_0$, with
\begin{align}
  \frac{d\mu}{d\mu_0}(\phi)
  &\;=\; Z^{-1}e^{-V(\phi)},
  \label{eq:mudensity}\\
  \frac{d(T_\psi)^*\mu}{d\mu_0}(\phi)
  &\;=\; Z^{-1}R_\psi(\phi)\, e^{-V(\phi-\psi)}.
  \label{eq:TmudensityGFF}
\end{align}
\end{lemma}

\begin{proof}
Equation~\eqref{eq:mudensity} is the definition of the $\Phi^4_3$ measure: $d\mu = Z^{-1}e^{-V}d\mu_0$.

For \eqref{eq:TmudensityGFF}, let $A\subset\mathcal{D}'(\mathbb{T}^3)$ be a Borel set.  By the
definition of pushforward:
\begin{equation}
  (T_\psi)^*\mu(A)
  \;=\; \mu\bigl(\{\phi : \phi + \psi \in A\}\bigr)
  \;=\; \int_{\{\phi:\phi+\psi\in A\}} Z^{-1}e^{-V(\phi)}\,d\mu_0(\phi).
  \label{eq:pushdef}
\end{equation}
We perform the change of variables by using the identity provided by Lemma~\ref{lem:CM}.
Since $\frac{d(T_\psi)^*\mu_0}{d\mu_0} = R_\psi$, for any non-negative measurable $g$:
\[
  \int g(\xi)\,d(T_\psi)^*\mu_0(\xi)
  \;=\; \int g(\xi)\, R_\psi(\xi)\,d\mu_0(\xi).
\]
On the left: $\int g(\xi)\,d(T_\psi)^*\mu_0(\xi) = \int g(\phi+\psi)\,d\mu_0(\phi)$.
Therefore:
\begin{equation}
  \int g(\phi+\psi)\,d\mu_0(\phi)
  \;=\; \int g(\xi)\, R_\psi(\xi)\,d\mu_0(\xi).
  \label{eq:CMformula}
\end{equation}
Substituting $g(\xi) = \mathbf{1}_A(\xi)\, e^{-V(\xi-\psi)}$ into \eqref{eq:pushdef}:
\[
  (T_\psi)^*\mu(A)
  \;=\; Z^{-1}\int g(\phi+\psi)\,d\mu_0(\phi)
  \;=\; Z^{-1}\int_A e^{-V(\xi-\psi)}\, R_\psi(\xi)\,d\mu_0(\xi),
\]
which establishes \eqref{eq:TmudensityGFF}.

We must check that $V(\xi-\psi)$ is well-defined $\mu_0$-a.s.  The $\Phi^4_3$ measure is
constructed so that $V(\phi)$ is a finite real number for $\mu_0$-a.e.\ $\phi$ (this is the
content of the renormalization procedure: the Wick-ordered monomials ${:\phi^k:}_{\mu_0}$
are well-defined as random variables in $L^2(\mu_0)$, cf.\ \cite[Section 4]{GubinelliHofmanova2019}).
Since $\psi\in C^\infty(\mathbb{T}^3)$ is a fixed smooth function, for $\mu_0$-a.e.\ $\phi$
the distribution $\phi - \psi$ is in the same regularity class as $\phi$, and
$V(\phi-\psi)$ is therefore finite $\mu_0$-a.s.\ by the same argument.
\end{proof}

\begin{lemma}[Positivity and $L^1$ integrability of the Radon--Nikodym derivative]
\label{lem:RN}
The function $G:\mathcal{D}'(\mathbb{T}^3)\to(0,\infty)$ defined by
\[
  G(\phi) \;:=\; R_\psi(\phi)\cdot e^{W(\phi-\psi;\,\psi)},
  \label{eq:Gdef}
\]
with $W(\xi;\psi) = V(\xi+\psi)-V(\xi)$, satisfies:
\begin{enumerate}
  \item[(i)] $G(\phi) > 0$ for $\mu$-a.e.\ $\phi$.
  \item[(ii)] $G \in L^1(\mu)$ with $\int G\,d\mu = 1$.
\end{enumerate}
\end{lemma}

\begin{proof}
\textit{Proof of (i).}
Since $R_\psi(\phi) = \exp(\langle(-\Delta+m^2)\psi,\phi\rangle - \frac{1}{2}\langle(-\Delta+m^2)\psi,\psi\rangle)$
is an exponential of a real-valued quantity, $R_\psi(\phi)\in(0,\infty)$ for every $\phi$.

The exponent $W(\phi-\psi;\psi) = V(\phi) - V(\phi-\psi)$ is finite $\mu$-a.s., because
$V(\phi)$ and $V(\phi-\psi)$ are both finite $\mu$-a.s.\ (as $\mu\ll\mu_0$ and they are
$\mu_0$-a.s.\ finite by Lemma~\ref{lem:densities}).  Therefore $e^{W(\phi-\psi;\psi)}\in(0,\infty)$
$\mu$-a.s.

Combining: $G(\phi) = R_\psi(\phi)\cdot e^{W(\phi-\psi;\psi)} > 0$ for $\mu$-a.e.\ $\phi$.

\textit{Proof of (ii).}
We compute directly using Lemmas~\ref{lem:CM} and \ref{lem:densities}:
\begin{align}
  \int G\,d\mu
  &\;=\; \int R_\psi(\phi)\, e^{V(\phi)-V(\phi-\psi)}\,d\mu(\phi) \nonumber \\
  &\;=\; \int R_\psi(\phi)\, e^{V(\phi)-V(\phi-\psi)}\cdot Z^{-1}e^{-V(\phi)}\,d\mu_0(\phi) \nonumber \\
  &\;=\; Z^{-1}\int R_\psi(\phi)\, e^{-V(\phi-\psi)}\,d\mu_0(\phi).
  \label{eq:intG}
\end{align}
By \eqref{eq:TmudensityGFF}, the right-hand side of \eqref{eq:intG} equals
$\int d(T_\psi)^*\mu = 1$ (since $(T_\psi)^*\mu$ is a probability measure).
Therefore $\int G\,d\mu = 1 < \infty$, so $G\in L^1(\mu)$.
\end{proof}

\begin{proof}[Proof of Theorem~\ref{thm:main}]
We establish the formula \eqref{eq:RNderiv} and then conclude equivalence.

\textit{Step 1: Formula for $d(T_\psi)^*\mu / d\mu$.}
By \eqref{eq:mudensity} and \eqref{eq:TmudensityGFF}, both $\mu$ and $(T_\psi)^*\mu$ are
absolutely continuous with respect to $\mu_0$.  Their Radon--Nikodym derivatives w.r.t.\ $\mu_0$
are, respectively, $Z^{-1}e^{-V(\phi)}$ and $Z^{-1}R_\psi(\phi)e^{-V(\phi-\psi)}$.  Therefore:
\begin{align}
  \frac{d(T_\psi)^*\mu}{d\mu}(\phi)
  &\;=\; \frac{d(T_\psi)^*\mu/d\mu_0(\phi)}{d\mu/d\mu_0(\phi)}
  \;=\; \frac{Z^{-1}R_\psi(\phi)\,e^{-V(\phi-\psi)}}{Z^{-1}e^{-V(\phi)}} \nonumber \\
  &\;=\; R_\psi(\phi)\cdot e^{V(\phi)-V(\phi-\psi)}
  \;=\; R_\psi(\phi)\cdot e^{W(\phi-\psi;\,\psi)},
  \label{eq:RNcomputed}
\end{align}
where the last equality uses the definition $W(\xi;\psi) = V(\xi+\psi)-V(\xi)$
(applied with $\xi = \phi-\psi$, so $V(\xi+\psi) - V(\xi) = V(\phi)-V(\phi-\psi)$).

The chain rule for Radon--Nikodym derivatives (applied to the $\sigma$-finite dominating
measure $\mu_0$) is valid precisely when both numerator and denominator are strictly positive
$\mu_0$-a.e.\ The denominator $Z^{-1}e^{-V(\phi)} > 0$ $\mu_0$-a.e.\ by the construction
of $\mu$.  The numerator $Z^{-1}R_\psi(\phi)e^{-V(\phi-\psi)} > 0$ $\mu_0$-a.e.\ because
$R_\psi>0$ everywhere (Lemma~\ref{lem:CM}) and $e^{-V(\phi-\psi)}>0$ $\mu_0$-a.e.

\textit{Step 2: $(T_\psi)^*\mu \ll \mu$.}
For any Borel set $A$ with $\mu(A) = 0$:
\[
  (T_\psi)^*\mu(A)
  \;=\; \int_A G\,d\mu
  \;=\; 0,
\]
where $G = d(T_\psi)^*\mu/d\mu$ as computed in \eqref{eq:RNcomputed}.  (The equality
$\int_A G\,d\mu = 0$ holds because $\mu(A)=0$ and $G$ is $\mu$-integrable by
Lemma~\ref{lem:RN}(ii).)  Hence $(T_\psi)^*\mu\ll\mu$.

\textit{Step 3: $\mu \ll (T_\psi)^*\mu$.}
We show: if $(T_\psi)^*\mu(A)=0$, then $\mu(A)=0$.  We have
\[
  0 \;=\; (T_\psi)^*\mu(A) \;=\; \int_A G\,d\mu,
\]
with $G(\phi)>0$ for $\mu$-a.e.\ $\phi$ (Lemma~\ref{lem:RN}(i)) and $G\in L^1(\mu)$
(Lemma~\ref{lem:RN}(ii)).  Since the integrand $G\mathbf{1}_A$ is non-negative and has
zero integral under $\mu$, it is zero $\mu$-a.e., i.e., $G(\phi) = 0$ for $\mu$-a.e.\ $\phi\in A$.
But $G>0$ $\mu$-a.e., so $\mu(A) = 0$.  Hence $\mu\ll (T_\psi)^*\mu$.

\textit{Conclusion.}
Steps~2 and 3 give $(T_\psi)^*\mu\sim\mu$.
\end{proof}

\begin{remark}
The answer to the question is \textbf{yes}: $\mu$ and $(T_\psi)^*\mu$ are equivalent.  The
key ingredients are (1) the Cameron--Martin theorem for the GFF, which requires only
$\psi\in H^1(\mathbb{T}^3)$ (satisfied since $\psi\in C^\infty$), and (2) the fact that both
the $\Phi^4_3$ density $e^{-V}$ and its shift $e^{-V(\cdot-\psi)}$ are simultaneously
$\mu_0$-a.s.\ finite and positive, which follows from the regularity of the $\Phi^4_3$
construction.  No quantitative information about the size of $\psi$ is needed; equivalence
holds for any nonzero $\psi\in C^\infty(\mathbb{T}^3)$.
\end{remark}

\begin{remark}
The Radon--Nikodym derivative \eqref{eq:RNderiv} can be decomposed as
\begin{align*}
  G(\phi)
  &= \exp\!\Bigl(\langle(-\Delta+m^2)\psi,\,\phi\rangle_{L^2}
      - \tfrac{1}{2}\langle(-\Delta+m^2)\psi,\,\psi\rangle_{L^2}\Bigr) \\
  &\quad\times\exp\!\Bigl(V(\phi) - V(\phi-\psi)\Bigr).
\end{align*}
The factor $V(\phi)-V(\phi-\psi)$ equals, when expanded using the Wick polynomial shift
identity $:(\xi+\psi)^k:_{\mu_0} = \sum_{j=0}^k \binom{k}{j}\psi^{k-j}:\xi^j:_{\mu_0}$,
a polynomial of degree $3$ in Wick powers of $\phi-\psi$ with smooth $\psi$-dependent
coefficients.  The precise integrability of $G$ in $L^p(\mu)$ for $p>1$ can be analyzed
using hypercontractivity (Nelson's theorem) for the chaos components of degree $\leq 3$,
but for the equivalence of measures we only need $G\in L^1(\mu)$ (proved above) and
$G>0$ $\mu$-a.s.\ (proved above).
\end{remark}

\begin{thebibliography}{9}

\bibitem{Bogachev1998}
V.~I.~Bogachev,
\emph{Gaussian Measures},
Mathematical Surveys and Monographs, vol.~62,
American Mathematical Society, Providence, RI, 1998.

\bibitem{GlimmJaffe1987}
J.~Glimm and A.~Jaffe,
\emph{Quantum Physics: A Functional Integral Point of View},
2nd ed., Springer-Verlag, New York, 1987.

\bibitem{GubinelliHofmanova2019}
M.~Gubinelli and M.~Hofmanov\'{a},
A PDE construction of the Euclidean $\Phi^4_3$ quantum field theory,
\emph{Commun.\ Math.\ Phys.}, \textbf{384} (2021), 1–75.

\bibitem{Hairer2014}
M.~Hairer,
A theory of regularity structures,
\emph{Invent.\ Math.}, \textbf{198} (2014), no.~2, 269–504.

\bibitem{Nelson1966}
E.~Nelson,
A quartic interaction in two dimensions,
in \emph{Mathematical Theory of Elementary Particles}, MIT Press, Cambridge, MA, 1966,
pp.~69–73.

\end{thebibliography}

\end{document}
